% Half title, title page, copyright page, dedication
\thispagestyle{empty}
\begin{center}
\vspace*{2.2in}
{\Huge\scshape On the Collatz Conjecture}
\end{center}
\cleardoublepage

\thispagestyle{empty}
\begin{center}
\vspace*{1.1in}
{\huge\scshape On the Collatz Conjecture}\\[16pt]
{\Large\itshape Martingales, Momentum\\[4pt] and Continued Fractions}\\[1.2in]
{\Large Nick Polson}\\[4pt]
{\small Booth School of Business, University of Chicago}\\[22pt]
{\Large Vadim Sokolov}\\[4pt]
{\small George Mason University}
\vfill
{\small 2026}
\end{center}
\clearpage

\thispagestyle{empty}
\vspace*{\fill}
{\small\noindent
Copyright \copyright\ 2026 Nicholas G.\ Polson and Vadim Sokolov.\\[8pt]
On the Collatz Conjecture: Martingales, Momentum and Continued Fractions.\\[8pt]
Mathematics Subject Classification (2020): 11B83, 11J70, 37P99, 60G42, 90C25.\\[8pt]
Every number computed in this book comes from code that accompanies it, and the scripts are listed in \cref{app:code}; numbers quoted from the literature are cited where they appear.
The source, the code and the digital edition are at \url{https://github.com/VadimSokolov/collatz}, and the book's page is \url{https://vsokolov.org/books/collatz/}.
The figures are in colour in the digital edition and in grayscale in the paperback.
The text was typeset in \LaTeX\ with
the Palatino typeface.\\[8pt]
First edition, October 2026.%
\ifx\BookISBNline\empty\else\\[8pt]\BookISBNline\fi
\ifx\BookDOIline\empty\else\\[8pt]\BookDOIline\fi
\\[8pt]License: distributed under the Creative Commons Attribution 4.0 International License (CC BY 4.0), \url{https://creativecommons.org/licenses/by/4.0/}: you may copy, share and adapt this book for any purpose, provided you credit the authors. The code is released under the MIT License.\\[8pt]
Independently published.\par}
\cleardoublepage

\thispagestyle{empty}
\vspace*{2in}
\begin{center}
\itshape
For the students who keep computing the orbit of $27$.
\end{center}
\cleardoublepage
